\uuid{5EIm}
\exo7id{3412}
\titre{exo7 3412}
\auteur{quercia}
\organisation{exo7}
\datecreate{2010-03-10}
\isIndication{false}
\isCorrection{true}
\chapitre{Déterminant, système linéaire}
\sousChapitre{Système linéaire, rang}
\cpge{mpsi-mp2i.calcul-matriciel-systemes.c, pcsi.calcul-matriciel-systemes-lineaires.c, ptsi.calcul-matriciel-systemes-lineaires.c, mpsi-mp2i.matrices-applications-lineaires.b, pcsi.matrices-applications-lineaires.b, ptsi.matrices-applications-lineaires.b, mpsi-mp2i.determinants.d, pcsi.calcul-matriciel-systemes-lineaires.d, ptsi.calcul-matriciel-systemes-lineaires.d, bcpst1.systemes-lineaires.a, bcpst1.systemes-lineaires.c, pt.determinants.a}
\module{Algèbre}
\niveau{L2}
\difficulte{}

\contenu{
\question{
\'Etudier l'existence de solutions du système :
$$\left\{\begin{array}{ccccccc}      x &+&       y &+&       z &= &1        \cr
                      ax &+&      by &+&      cz &= &d        \cr
                 a(a-1)x &+& b(b-1)y &+& c(c-1)z &= &d(d-1).  \cr\end{array}\right.$$
}
\reponse{
Système de Cramer ssi $a,b,c$ sont distincts. Sinon, il y a des solutions ssi $d\in\{a,b,c\}$.
}
}
